%%%%%%%%%%%%%%%%%%%%%%%%%%%%%%%%%%%%%%%%%%%%%%%%%%%%%%%%%%%%%%%%%%%%%%%%%%%%%%%%
% Frontiers in Systems Neuroscience -- Research Topic 72013
% "Embodied Interfaces: Human Experience in Virtual and Mediated Worlds"
%
% Extension of the MCC thesis (RIDL / LGAF-CNN / CFS) with re-run, corrected
% analyses. Every number in Results comes from results_macros.tex, generated by
% ../experiments/make_macros.py from the experiment outputs.
%%%%%%%%%%%%%%%%%%%%%%%%%%%%%%%%%%%%%%%%%%%%%%%%%%%%%%%%%%%%%%%%%%%%%%%%%%%%%%%%

\documentclass[utf8]{frontiers_suppmat}
\usepackage{multicol}
\usepackage[T1]{fontenc}
\usepackage{url,lineno,microtype}
\usepackage[onehalfspacing]{setspace}
\usepackage{amsmath,amssymb,amsfonts,amsthm,mathtools,bm}
\usepackage{hyperref}
\usepackage{graphicx,booktabs,multirow,siunitx}
\usepackage{xcolor}
\newcommand{\TODO}[1]{\textcolor{red}{[TODO: #1]}}
% frontiers_suppmat.cls has no \keywords (it comes from the main Frontiers template)
\providecommand{\keywords}[1]{\par\medskip\noindent\textbf{Keywords:} #1}
\graphicspath{{figures/}{../figures/}{../frontiers_ext/figures/}}
% AUTO-GENERATED by experiments/make_macros.py -- do not edit by hand
\newcommand{\hCfsVR}{2.21}
\newcommand{\hCfsVRsd}{0.25}
\newcommand{\hCfsNVR}{2.10}
\newcommand{\hCfsNVRsd}{0.28}
\newcommand{\hCfsT}{1.02}
\newcommand{\hCfsP}{=0.33}
\newcommand{\hCfsPholm}{=1.00}
\newcommand{\hCfsDz}{0.32}
\newcommand{\hCfsBF}{2.11}
\newcommand{\hCfsCIlo}{-0.08}
\newcommand{\hCfsCIhi}{0.30}
\newcommand{\hCfsTost}{=0.17}
\newcommand{\hCfsTostDz}{=0.30}
\newcommand{\hCfsNvr}{4}
\newcommand{\hCfsCvVR}{0.40}
\newcommand{\hCfsCvVRsd}{0.05}
\newcommand{\hCfsCvNVR}{0.76}
\newcommand{\hCfsCvNVRsd}{0.57}
\newcommand{\hCfsCvT}{-2.12}
\newcommand{\hCfsCvP}{=0.06}
\newcommand{\hCfsCvPholm}{=0.32}
\newcommand{\hCfsCvDz}{-0.67}
\newcommand{\hCfsCvBF}{0.68}
\newcommand{\hCfsCvCIlo}{-0.67}
\newcommand{\hCfsCvCIhi}{-0.05}
\newcommand{\hCfsCvTost}{=0.94}
\newcommand{\hCfsCvTostDz}{=0.70}
\newcommand{\hCfsCvNvr}{4}
\newcommand{\hOptVR}{0.939}
\newcommand{\hOptVRsd}{0.024}
\newcommand{\hOptNVR}{0.932}
\newcommand{\hOptNVRsd}{0.028}
\newcommand{\hOptT}{0.60}
\newcommand{\hOptP}{=0.56}
\newcommand{\hOptPholm}{=1.00}
\newcommand{\hOptDz}{0.19}
\newcommand{\hOptBF}{2.78}
\newcommand{\hOptCIlo}{-0.013}
\newcommand{\hOptCIhi}{0.027}
\newcommand{\hOptTost}{<0.001}
\newcommand{\hOptTostDz}{=0.18}
\newcommand{\hOptNvr}{3}
\newcommand{\hLdljVR}{-20.25}
\newcommand{\hLdljVRsd}{0.39}
\newcommand{\hLdljNVR}{-20.30}
\newcommand{\hLdljNVRsd}{0.57}
\newcommand{\hLdljT}{0.23}
\newcommand{\hLdljP}{=0.82}
\newcommand{\hLdljPholm}{=1.00}
\newcommand{\hLdljDz}{0.07}
\newcommand{\hLdljBF}{3.17}
\newcommand{\hLdljCIlo}{-0.38}
\newcommand{\hLdljCIhi}{0.49}
\newcommand{\hLdljTost}{<0.001}
\newcommand{\hLdljTostDz}{=0.10}
\newcommand{\hLdljNvr}{7}
\newcommand{\hMicroVR}{0.97}
\newcommand{\hMicroVRsd}{0.67}
\newcommand{\hMicroNVR}{1.26}
\newcommand{\hMicroNVRsd}{1.11}
\newcommand{\hMicroT}{-0.70}
\newcommand{\hMicroP}{=0.50}
\newcommand{\hMicroPholm}{=1.00}
\newcommand{\hMicroDz}{-0.22}
\newcommand{\hMicroBF}{2.64}
\newcommand{\hMicroCIlo}{-1.03}
\newcommand{\hMicroCIhi}{0.46}
\newcommand{\hMicroTost}{=0.65}
\newcommand{\hMicroTostDz}{=0.20}
\newcommand{\hMicroNvr}{6}
\newcommand{\hPathEffVR}{0.36}
\newcommand{\hPathEffVRsd}{0.16}
\newcommand{\hPathEffNVR}{0.26}
\newcommand{\hPathEffNVRsd}{0.13}
\newcommand{\hPathEffT}{2.40}
\newcommand{\hPathEffP}{=0.04}
\newcommand{\hPathEffPholm}{=0.40}
\newcommand{\hPathEffDz}{0.76}
\newcommand{\hPathEffBF}{0.48}
\newcommand{\hPathEffCIlo}{0.02}
\newcommand{\hPathEffCIhi}{0.17}
\newcommand{\hPathEffTost}{=0.94}
\newcommand{\hPathEffTostDz}{=0.78}
\newcommand{\hPathEffNvr}{8}
\newcommand{\hPathLenVR}{0.29}
\newcommand{\hPathLenVRsd}{0.16}
\newcommand{\hPathLenNVR}{0.36}
\newcommand{\hPathLenNVRsd}{0.19}
\newcommand{\hPathLenT}{-0.89}
\newcommand{\hPathLenP}{=0.40}
\newcommand{\hPathLenPholm}{=1.00}
\newcommand{\hPathLenDz}{-0.28}
\newcommand{\hPathLenBF}{2.34}
\newcommand{\hPathLenCIlo}{-0.20}
\newcommand{\hPathLenCIhi}{0.07}
\newcommand{\hPathLenTost}{=0.65}
\newcommand{\hPathLenTostDz}{=0.25}
\newcommand{\hPathLenNvr}{4}
\newcommand{\hActiveVR}{0.69}
\newcommand{\hActiveVRsd}{0.08}
\newcommand{\hActiveNVR}{0.72}
\newcommand{\hActiveNVRsd}{0.09}
\newcommand{\hActiveT}{-0.75}
\newcommand{\hActiveP}{=0.47}
\newcommand{\hActivePholm}{=1.00}
\newcommand{\hActiveDz}{-0.24}
\newcommand{\hActiveBF}{2.56}
\newcommand{\hActiveCIlo}{-0.10}
\newcommand{\hActiveCIhi}{0.04}
\newcommand{\hActiveTost}{=0.14}
\newcommand{\hActiveTostDz}{=0.21}
\newcommand{\hActiveNvr}{5}
\newcommand{\hSubmVR}{9.1}
\newcommand{\hSubmVRsd}{3.5}
\newcommand{\hSubmNVR}{7.5}
\newcommand{\hSubmNVRsd}{2.8}
\newcommand{\hSubmT}{1.48}
\newcommand{\hSubmP}{=0.17}
\newcommand{\hSubmPholm}{=1.00}
\newcommand{\hSubmDz}{0.47}
\newcommand{\hSubmBF}{1.39}
\newcommand{\hSubmCIlo}{-0.4}
\newcommand{\hSubmCIhi}{3.6}
\newcommand{\hSubmTost}{=0.77}
\newcommand{\hSubmTostDz}{=0.46}
\newcommand{\hSubmNvr}{7}
\newcommand{\hCfsCIloPct}{-4}
\newcommand{\hCfsCIhiPct}{14}
\newcommand{\mmBeta}{0.11}
\newcommand{\mmSE}{0.041}
\newcommand{\mmP}{=0.007}
\newcommand{\mmN}{598}
\newcommand{\mmInterP}{=0.50}
\newcommand{\mmRsBeta}{0.10}
\newcommand{\mmRsSE}{0.064}
\newcommand{\mmRsP}{=0.12}
\newcommand{\mmTrimBeta}{0.08}
\newcommand{\mmTrimSE}{0.081}
\newcommand{\mmTrimP}{=0.30}
\newcommand{\mmRankBeta}{0.12}
\newcommand{\mmRankSE}{0.071}
\newcommand{\mmRankP}{=0.10}
\newcommand{\mmLdljBeta}{-0.02}
\newcommand{\mmLdljP}{=0.61}
\newcommand{\mmMicroBeta}{-0.06}
\newcommand{\mmMicroP}{=0.17}
\newcommand{\msNup}{8}
\newcommand{\msHi}{0.941}
\newcommand{\msLo}{0.931}
\newcommand{\msP}{=0.11}
\newcommand{\msDz}{0.57}
\newcommand{\aRidlErr}{13.79}
\newcommand{\aRidlCi}{3.18}
\newcommand{\aRidlEtraj}{13.57}
\newcommand{\aRidlSig}{2.13}
\newcommand{\aRidlNmse}{4.359}
\newcommand{\aRidlSteps}{14,826}
\newcommand{\aRidlN}{20}
\newcommand{\aPpoErr}{15.33}
\newcommand{\aPpoCi}{1.93}
\newcommand{\aPpoEtraj}{15.17}
\newcommand{\aPpoSig}{1.96}
\newcommand{\aPpoNmse}{4.050}
\newcommand{\aPpoSteps}{14,826}
\newcommand{\aPpoN}{20}
\newcommand{\aBcErr}{10.68}
\newcommand{\aBcCi}{3.92}
\newcommand{\aBcEtraj}{10.38}
\newcommand{\aBcSig}{2.20}
\newcommand{\aBcNmse}{3.851}
\newcommand{\aBcSteps}{0}
\newcommand{\aBcN}{20}
\newcommand{\aNoVisErr}{8.42}
\newcommand{\aNoVisCi}{1.14}
\newcommand{\aNoVisEtraj}{8.20}
\newcommand{\aNoVisSig}{2.03}
\newcommand{\aNoVisNmse}{2.457}
\newcommand{\aNoVisSteps}{14,826}
\newcommand{\aNoVisN}{20}
\newcommand{\aPixErr}{15.00}
\newcommand{\aPixCi}{2.00}
\newcommand{\aPixEtraj}{14.83}
\newcommand{\aPixSig}{2.50}
\newcommand{\aPixNmse}{3.758}
\newcommand{\aPixSteps}{14,826}
\newcommand{\aPixN}{20}
\newcommand{\aNoFlowRErr}{13.68}
\newcommand{\aNoFlowRCi}{1.53}
\newcommand{\aNoFlowREtraj}{13.46}
\newcommand{\aNoFlowRSig}{2.08}
\newcommand{\aNoFlowRNmse}{4.322}
\newcommand{\aNoFlowRSteps}{14,826}
\newcommand{\aNoFlowRN}{20}
\newcommand{\aNoCfsInErr}{13.63}
\newcommand{\aNoCfsInCi}{3.44}
\newcommand{\aNoCfsInEtraj}{13.44}
\newcommand{\aNoCfsInSig}{2.06}
\newcommand{\aNoCfsInNmse}{4.561}
\newcommand{\aNoCfsInSteps}{14,826}
\newcommand{\aNoCfsInN}{20}
\newcommand{\aPooledErr}{14.05}
\newcommand{\aPooledCi}{1.97}
\newcommand{\aPooledEtraj}{13.76}
\newcommand{\aPooledSig}{2.15}
\newcommand{\aPooledNmse}{3.678}
\newcommand{\aPooledSteps}{14,865}
\newcommand{\aPooledN}{20}
\newcommand{\aHoldErr}{4.93}
\newcommand{\aHoldCi}{0.98}
\newcommand{\aHoldEtraj}{4.54}
\newcommand{\aHoldSig}{2.37}
\newcommand{\aHoldNmse}{2.093}
\newcommand{\aHoldSteps}{0}
\newcommand{\aHoldN}{10}
\newcommand{\aCVErr}{11.85}
\newcommand{\aCVCi}{2.57}
\newcommand{\aCVEtraj}{11.15}
\newcommand{\aCVSig}{6.73}
\newcommand{\aCVNmse}{152.087}
\newcommand{\aCVSteps}{0}
\newcommand{\aCVN}{10}
\newcommand{\dPpo}{+1.54}
\newcommand{\dPpoP}{=0.26}
\newcommand{\dBc}{-3.11}
\newcommand{\dBcP}{=0.001}
\newcommand{\dNoVis}{-5.37}
\newcommand{\dNoVisP}{<0.001}
\newcommand{\dPix}{+1.22}
\newcommand{\dPixP}{=0.11}
\newcommand{\dNoFlowR}{-0.11}
\newcommand{\dNoFlowRP}{=0.45}
\newcommand{\dNoCfsIn}{-0.16}
\newcommand{\dNoCfsInP}{=0.52}
\newcommand{\dPooled}{+0.26}
\newcommand{\dPooledP}{=0.43}
\newcommand{\dHold}{-8.85}
\newcommand{\dHoldP}{=0.002}
\newcommand{\dCV}{-1.93}
\newcommand{\dCVP}{=0.77}
\newcommand{\srcVR}{12.19}
\newcommand{\srcNVR}{15.39}
\newcommand{\srcDiff}{-3.20}
\newcommand{\srcP}{=0.35}
\newcommand{\srcTost}{=0.69}
\newcommand{\srcCIlo}{-9.13}
\newcommand{\srcCIhi}{+2.72}
\newcommand{\cRidlFinal}{5.080}
\newcommand{\cPpoFinal}{6.392}
\newcommand{\cRidlFirst}{5.370}
\newcommand{\cSteps}{36,900}
\newcommand{\cPpoReach}{100}
\newcommand{\cPpoFirst}{3.730}
\newcommand{\tBcErr}{99.1}
\newcommand{\tBcNmse}{144.76}
\newcommand{\tRidlErr}{219.6}
\newcommand{\tRidlNmse}{230.74}
\newcommand{\tBase}{22.6}
\newcommand{\lgafMaxDz}{0.26}
\newcommand{\lgafMinP}{0.44}
\newcommand{\lgafWin}{23}
\newcommand{\lgafWinTotal}{466}

% frontiers_suppmat.cls numbers everything S1, S2, ... -- this is the main text
\renewcommand\theequation{\arabic{equation}}
\renewcommand\thefigure{\arabic{figure}}
\renewcommand\thetable{\arabic{table}}
% figure that is drawn if present, otherwise a visible placeholder (pending results)
\newcommand{\figfile}[2]{\IfFileExists{figures/#1}{\includegraphics[width=#2\textwidth]{#1}}{%
  \IfFileExists{../figures/#1}{\includegraphics[width=#2\textwidth]{#1}}{%
  \fbox{\parbox[c][3.5cm][c]{#2\textwidth}{\centering\TODO{\texttt{\detokenize{#1}} is generated by run\_all.sh}}}}}}

% Layout: no vertical stretching to fill pages, and let floats share pages with text
\raggedbottom
\renewcommand{\topfraction}{0.9}
\renewcommand{\bottomfraction}{0.8}
\renewcommand{\textfraction}{0.07}
\renewcommand{\floatpagefraction}{0.85}
\setcounter{topnumber}{3}
\setcounter{bottomnumber}{2}
\setcounter{totalnumber}{4}
\predisplaypenalty=200 \clubpenalty=150 \widowpenalty=150 \displaywidowpenalty=150
\begin{document}
\onecolumn
\firstpage{1}

\title[Neuroadaptive RIDL: Cognitive Equivalence of VR--Haptic Teleoperation]{Neuroadaptive Decision-Making AI for Immersive Surgical Robotics Training: Cognitive Equivalence of VR--Haptic Teleoperation and a Hierarchical Learning Agent Ablation}

\maketitle
% frontiers_suppmat.cls does not typeset \author/\address, so the block is set by hand.
{\fontsize{11}{15}\selectfont\raggedright
Nezih Nieto Guti\'errez$^{1}$, Alejandro Hern\'andez G\'omez$^{1}$, Gustavo De Los R\'ios Alatorre$^{1}$, Felipe de Jes\'us F\'elix Arredondo$^{3}$, Arturo Jos\'e Murra L\'opez$^{2}$ and Luis Alberto Mu\~noz Ubando$^{1,*}$\par}
\smallskip
{\small\raggedright
$^{1}$School of Engineering and Sciences, Tecnol\'ogico de Monterrey, Monterrey, Mexico\par
$^{2}$University of Colorado Anschutz Medical Campus, Denver, CO, USA\par
$^{3}$Department of Research, Healthen Group, Monterrey, NL, Mexico\par
$^{*}$Correspondence: Luis Alberto Mu\~noz Ubando, \url{amunoz@tec.mx}\par}
\medskip

% -----------------------------------------------------------------
\section*{Abstract}
\begin{multicols}{2}
Surgical robotics demands sustained attention, stable motor coordination and rapid decision-making, yet trainees are still assessed mostly by what they achieve rather than by how they engage while achieving it. Immersive interfaces that combine virtual reality (VR) with haptic feedback are often assumed to help learners, but a more basic requirement precedes any such benefit: the interface must not \emph{cost} the learner cognitively, since a neuroadaptive tutor would otherwise read an engagement signal distorted by the very medium through which training is delivered. This study tested that requirement within a neuroadaptive training framework that integrates VR, force feedback, surgical teleoperation simulation and wearable electroencephalography (EEG). Ten participants teleoperated a simulated da Vinci Research Kit (dVRK) through a Meta Quest~3 headset with Novint Falcon force feedback and, in a second session, through a 2D monitor without haptics, while 8-channel dry EEG, tool kinematics and endoscopic video were recorded. Engagement was quantified by a continuous Cognitive Flow Score (CFS), obtained from a learnable Gramian Angular Field convolutional network (LGAF-CNN) and scaled by questionnaire-virtualized $\beta/\alpha$ power. Immersion did not reduce engagement (CFS \hCfsVR{} vs.\ \hCfsNVR{}; $p\hCfsP$, BF$_{01}=\hCfsBF$), and jerk-based optimality and smoothness were statistically equivalent between interfaces within $\pm$10\% (TOST $p\hOptTost$). Within participants, moments of higher CFS coincided with more optimal movement ($\beta=\mmBeta$, $p\mmP$), although the strength of this coupling differed between individuals. The same recordings then served as demonstrations for RIDL, a hierarchical learning agent that pre-trains a visuomotor policy built on a frozen Stable-Diffusion encoder by imitation and refines it with Proximal Policy Optimization (PPO) on a CFS-weighted reward. In a subject-wise cross-validated ablation with 20 runs per variant, RIDL showed a balanced error and volatility profile, with lower mean error than PPO trained from scratch (\aRidlErr{} vs.\ \aPpoErr{}\,mm) and far more stable predictions than constant-velocity extrapolation, while lower-capacity variants and a hold-position reference remained more accurate, as expected from the 16 demonstrations available per fold. Taken together, the findings indicate that immersive VR--haptic teleoperation is cognitively transparent and can serve as the substrate for neuroadaptive training, and that demonstrations from either interface can train hierarchical agents whose full benefit awaits larger demonstration sets and closed-loop simulation.

\keywords{EEG, cognitive engagement, embodiment, surgical robotics, virtual reality, haptic feedback, equivalence testing, imitation learning, reinforcement learning, diffusion policy}
\end{multicols}

% =================================================================
\section{Introduction}

Robotic-assisted surgery has changed what a surgeon must learn, because beyond anatomy and procedure a trainee at a teleoperated console must perceive depth from an endoscopic view, translate hand motion into instrument motion through a kinematic mapping, regulate the forces applied to tissue and keep track of the whole surgical field while doing so \cite{Haidegger2022,Riveiro-Moreno2023,Peters2018,Fujie2020,Fiorini2022}. These demands are cognitive as much as they are motor, and they must be sustained for long periods. Training curricula, however, still judge competence largely after the fact, through task outcomes, completion times and expert ratings, and they say little about the learner's engagement while the skill is being acquired \cite{Attanasio2021,Vasey2023}. A trainee who completes a task while barely attending and one who completes it while fully engaged may receive the same score, even though they are on very different learning trajectories.

This gap can now be closed through two converging developments, since wearable biosensors, including electroencephalography (EEG), eye tracking and heart-rate variability, now give continuous access to attentional and affective states outside the laboratory \cite{10.3389/fnrgo.2022.1052411,Shafiei2023,Alcaide2021Focus}, and immersive simulation, which combines virtual reality (VR) with haptic rendering, provides safe and high-fidelity practice \cite{Munawar2023,Munawar2021,Alaker2016,Rangarajan2020}. Immersive simulation has been associated with improved skill acquisition \cite{Mao2021,Gawecki2020,9262701}, and haptic feedback increases simulator fidelity, particularly for novices \cite{hapticshci2,brugada2022}. When the two technologies are combined, the learner's state can be read during practice and fed back into the simulation, which is the premise of \emph{closed-loop}, neuroadaptive instruction \cite{Guger2009Brain,long2022,qian2019}.

The same richness that makes immersion attractive also gives reason for caution, since cognitive-load theory predicts that additional sensory channels can impose extraneous load that competes with the task itself \cite{jining2021lost,Sokolowska2023-cz}, and a stereoscopic headset with force feedback is a considerably richer environment than a monitor. For a neuroadaptive tutor this possibility is far from a side issue, because if the immersive interface altered the learner's engagement, the tutor's input signal would depend on the display through which training is delivered, and a policy calibrated in one setting could not be trusted in the other. The first question an embodied training interface must answer is therefore not whether it helps, but whether it is \emph{cognitively transparent}, that is, whether it leaves engagement and movement quality as they would be on a conventional display. Answering this question requires a different statistical stance from the usual one, because a non-significant difference does not demonstrate the absence of a difference; equivalence tests and Bayes factors are needed to provide positive evidence that no meaningful cost exists.

A second motivation concerns the agents on which such a tutor might eventually rely, and for which reinforcement learning (RL) offers an attractive route to surgical autonomy and assistance \cite{Zhang2021,Brunke2022,Kassahun2016,Andras2020}, but its sample inefficiency limits deployment in a domain where every interaction is costly \cite{Barnoy2021,Xu2021}. Learning from demonstrations offers a way around this limitation, and diffusion-based visuomotor policies and surgical robot transformers have shown that rich policies can be learned from recorded expert motion \cite{chi2023diffusion,kim2024srt}. Demonstrations recorded in a neuroadaptive simulator carry something that ordinary demonstrations do not: a continuous measure of how engaged the demonstrator was at each moment. Whether this measure can shape what an agent learns, for example by giving more weight to moments of high engagement, is an open question, and its practical value depends in turn on whether demonstrations recorded through immersive and conventional interfaces are interchangeable.

This study addresses both questions within a neuroadaptive training framework previously developed by the authors. The framework couples VR--haptic teleoperation of a da Vinci Research Kit (dVRK) simulator with an EEG-based Cognitive Flow Score (CFS), computed by a learnable Gramian Angular Field convolutional network (LGAF-CNN), and with RIDL, a hierarchical agent that learns surgical tool motion from the recorded demonstrations. Every component is described in full in the Materials and Methods, so that the study can be understood and reproduced without recourse to earlier reports.

The first of the study's three contributions is a within-subject test of cognitive and motor equivalence between VR+haptic and 2D teleoperation, in which equivalence tests and Bayes factors complement conventional hypothesis tests. The second is a moment-to-moment analysis of whether CFS tracks execution quality within each session, which asks whether the engagement signal is behaviorally meaningful and not merely stable across interfaces. The third is a subject-wise cross-validated ablation of RIDL, in which each hierarchical component (the imitation prior, the PPO refinement, the Stable-Diffusion visual prior, the CFS-weighted reward, the CFS state input and the subject curriculum) is removed in turn and the outcome is compared with non-learning reference predictors, together with a comparison of agents trained on VR and on 2D demonstrations. The article first describes the apparatus, measures and learning agent, then reports the human and agent results in that order, and closes by discussing what they imply for neuroadaptive surgical training.

% =================================================================
\section{Materials and Methods}
Testing whether an immersive interface preserves the learner's engagement requires three ingredients: a task that can be performed identically with and without immersion, a continuous measure of engagement that can be read from the brain while the learner moves, and objective measures of how well the tool is moved. The study obtained all three from the same sessions by combining a simulated surgical robot, a stereoscopic headset with force feedback, wearable EEG and synchronized tool logging. The account below follows the data from their origin to their use, beginning with the participants, the simulator and the protocol, continuing with the EEG recording, the circuit model that guides its interpretation and the construction of the engagement score, then defining the movement metrics and the statistical tests applied to them, and closing with the learning agent that reuses the recordings as demonstrations, so that every quantity reported in the Results is derived explicitly along the way.

\subsection{Participants}
Ten volunteers with novice-to-intermediate experience in robotic manipulation took part (age range 18--25 years; nine right-handed, one left-handed). This is the population for which training interfaces matter most and for which haptic feedback has been reported to be most beneficial \cite{hapticshci2,brugada2022}. Participants were required to have normal or corrected-to-normal vision and to tolerate VR and haptic devices, and neurological disorders, severe simulator sickness and motor impairments were grounds for exclusion. All participants gave written informed consent before taking part, and all data were anonymized before analysis.

\subsection{Training Environment}
A fair comparison requires that both interfaces present exactly the same surgical scene and tasks and differ only in how the scene is seen and whether contact is felt. A single simulator therefore drove both conditions, and only the display and the force channel were switched between them (Figure~\ref{fig:setup}).

\begin{figure}[!htb]
\centering
\figfile{fig1_setup.pdf}{0.95}
\caption{\textbf{Immersive neuroadaptive training setup.} (A) Participant wearing the Meta Quest~3 headset and Enobio EEG while teleoperating simulated dVRK instruments through a Novint Falcon. (B) System architecture: ROS orchestration, Docker-based haptics, GStreamer stereo streaming, OBS virtual camera, Unity/OpenXR rendering and synchronized logging of trajectories, forces, EEG and images. (C) EEG montage: the eight recorded sites (O1, Oz, O2, CP1, CP2, C3, C4, Fz) are highlighted.}
\label{fig:setup}
\end{figure}

The surgical field was rendered by a teleoperation simulator of the dVRK, a standard research platform for surgical robotics \cite{d2021accelerating,surrol,9807505}, built on AMBF \cite{8968568,Munawar2021} and reproducing the depth cues, tool articulation and endoscope motion typical of minimally invasive procedures \cite{9697399,9925319}. For the immersive condition, the stereoscopic endoscope view was streamed to a Meta Quest~3 headset through GStreamer and an OBS virtual camera and displayed with Unity and OpenXR. Throughout every session, the tool pose (position together with roll, pitch and yaw), the interaction forces and the left-endoscope frames were logged through ROS \cite{mower2022rospybullet}, which provided the common timeline on which all later signals were aligned.

Force feedback was delivered by a three-degree-of-freedom Novint Falcon \cite{hapticshci2} whose driver ran in a Docker container and was bridged to ROS by custom nodes, so that contact events in the simulation produced a proportional resistance at the participant's hand \cite{Rangarajan2020,Alaker2016}, a channel that was disabled in the 2D condition.

\subsection{Experimental Conditions and Protocol}
Every participant operated the simulated robot in two ways: in the VR+haptics condition, the stereoscopic scene was viewed through the headset and contact forces were rendered by the force-feedback device, whereas in the 2D/no-haptics condition (NVR) the same scene was viewed on a conventional monitor without force feedback. Task objectives, scoring and timing were identical in both, so that any difference between the conditions can be attributed to the display and to the presence of force feedback, subject to the fixed order of the sessions discussed at the end of this subsection.

All tasks took place in the same simulated suturing scene: a soft-tissue phantom with a curved incision whose entry and exit points were marked in red, a curved suture needle with its thread, and the patient-side instruments of the dVRK, viewed through the simulated endoscope. Each interface session followed the same sequence of four EEG-recorded blocks, which progressed from unconstrained familiarization to the most demanding manipulation. The session opened with a Reference block (2\,min of EEG), in which the participant moved a single instrument freely through the workspace without handling any object; this block familiarized the participant with the mapping between the haptic device and the tool and provided a baseline of unconstrained movement. In Task~1 (5\,min of EEG), the participant performed a pick-and-place of the needle, bringing the instrument to the needle, grasping it, carrying it across the phantom and releasing it at a new location, a grasp-and-transfer primitive of robotic surgical training \cite{Jackson2013,ManojVarier2020}. A resting-state block (3\,min of EEG) followed, during which the participant sat still without operating the device, which separated the two tasks and allowed engagement to return toward baseline. The session closed with Task~2 (5\,min of EEG), in which the participant attempted a suture: the needle was grasped and its tip driven across the marked entry and exit points of the incision while the thread followed, the most dexterous manipulation of the protocol \cite{Lobbezoo2021,Chiu2020}.

Because not every minute of EEG was accompanied by synchronized tool logging and video, the analyses rely throughout on the synchronized segments, which covered about 31\,s of the Reference block and about 77\,s of each task. The resting-state block contributed no tool data; it served to separate the tasks and as the resting condition for the questionnaires described in the next subsection. All participants performed the VR+haptics session first and the 2D/no-haptics session second, so the condition order was \emph{not} counterbalanced, and the consequences of this choice are considered in the Limitations.

\subsection{EEG Acquisition}
EEG was recorded throughout every block with an 8-channel Enobio dry-electrode system (NIC~2.1) at 500\,Hz, from O1, Oz, O2, CP1, CP2, C3, C4 and Fz, with a 50\,Hz line filter and the on-board accelerometer enabled \cite{Duvinage2013Performance,Qiu2019Assessing}. Dry consumer-grade systems trade some signal quality for rapid setup and comfort, which matters when participants also wear a headset, and they have been shown to be adequate for engagement indices when configured carefully \cite{Frey2016Comparison,Han2018Interactive}. Alongside the EEG, participants completed two questionnaires: the UWIST Mood Adjective Checklist, which scores energetic arousal, tense arousal and hedonic tone, and the Amsterdam Resting-State Questionnaire (ARSQ), which characterizes cognition during rest \cite{10.3389/fnrgo.2022.1052411}, and both scores enter the engagement score described below, while the choice of electrode sites follows from the circuit model introduced next.

\subsection{Decision-Making Circuit Model}
Reading engagement from a handful of scalp electrodes requires a view of which brain systems take part in surgical decisions and which of them the electrodes can reach, and the framework provides that view by treating surgical decision-making as a distributed circuit (Figure~\ref{fig:circuit}A) whose centre is a core decision network (Fz, FCz), coupled to goal maintenance (Fz), executive control (FC1/FC2) and motor planning and output (C3/C4). Reward processing (FCz, Fp1/Fp2) interacts with the basal ganglia and the subthalamic nucleus to weigh the value of an action against its cost, while dopaminergic (VTA) and noradrenergic (LC) centres modulate the balance between exploiting a known strategy and exploring alternatives. Attention enters through moderators in dorsolateral and ventromedial prefrontal cortex (F3/F4, Fp1/Fp2) and in parietal cortex (P3/P4, Pz).

The model shapes the rest of the study in two ways, the first of which concerns learning: in the model, the prefrontal--basal-ganglia--thalamic loop learns through dopamine-based reinforcement, which is what motivates the actor--critic structure of RIDL (Figure~\ref{fig:circuit}C); the correspondence is used as a design guide and is not offered as a claim about mechanism. The second concerns measurement, since the 8-channel montage samples the core decision and maintenance node (Fz), motor output (C3, C4), parietal attention (CP1 and CP2, the recorded sites nearest P3/P4) and visual input (O1, Oz, O2) (Figure~\ref{fig:circuit}B), whereas the deep structures of the circuit are not accessible to scalp EEG and can only be inferred indirectly.

\begin{figure}[!htb]
\centering
\figfile{fig_decision_circuit.pdf}{1.0}
\caption{\textbf{Decision-making circuit model.} (A) Compact-block model of the circuit with electrode loci (the circuit model underlying the framework): core decision (red), goal maintenance (blue), executive control (violet), motor planning and output (green), reward processing (yellow), basal ganglia/STN, and modulatory VTA and LC (orange); attentional moderators (dlPFC, vmPFC, parietal) in the top group. Solid lines: core interactions; dotted: modulatory projections; dashed: moderator input. Thick borders mark nodes sampled in this study. (B) Node loci on the 10--20 scalp layout; bold rings mark the eight recorded Enobio sites. (C) Correspondence used to design RIDL (analogy, not mechanism): prefrontal decision and control $\rightarrow$ policy head, reward/basal ganglia $\rightarrow$ value head, dopaminergic prediction error $\rightarrow$ PPO advantage, noradrenergic exploit/explore $\rightarrow$ policy variance and entropy bonus, attentional moderators $\rightarrow$ CFS as state input and reward weight, motor output $\rightarrow$ action chunk.}
\label{fig:circuit}
\end{figure}

\subsection{Cognitive Flow Score (CFS)}
With the recorded sites in place, engagement must still be turned into a quantity that varies from moment to moment, and the Cognitive Flow Score (CFS) serves this purpose, estimating task-directed engagement \cite{ROVIRA201789} in a form that is used both in the human analyses and as a learning signal for RIDL. Its construction starts from a classical spectral engagement index, replaces the part of that index that cannot be measured reliably during movement, and adds an EEG decoder that supplies the temporal dynamics.

\paragraph*{Spectral engagement index.} For an EEG channel $x(t)$, the spectrum is
\begin{equation}\label{eq:fourier}
X(f)=\int_{-\infty}^{\infty}x(t)\,e^{-j2\pi ft}\,dt ,
\end{equation}
and the log power in a canonical band is estimated with Welch's method,
\begin{equation}\label{eq:bandpower}
P_{\mathrm{band}}=\log\sum_{f\in\mathrm{band}}|X(f)|^{2},\qquad \mathrm{band}\in\{\delta,\theta,\alpha,\beta,\gamma\},
\end{equation}
with $\alpha$ = 8--13\,Hz and $\beta$ = 13--30\,Hz. The engagement index used in biocybernetic systems is the ratio
\begin{equation}\label{eq:engagement}
E=\frac{P_\beta}{P_\alpha},
\end{equation}
which increases when cortex shifts from idling ($\alpha$) to active processing ($\beta$).

\paragraph*{Virtualized band power.} During active manipulation, however, the spectrum is contaminated by motion and muscle artefacts, and a direct estimate of $P_\beta/P_\alpha$ would track movement as much as engagement. Band power is therefore \emph{virtualized} from the questionnaires administered with the EEG,
\begin{equation}\label{eq:virtual}
V_\beta\equiv\text{UWIST energetic arousal},\qquad V_\alpha\equiv\text{ARSQ total},\qquad E_v=\frac{V_\beta}{V_\alpha},
\end{equation}
where both scores are normalized before the ratio is taken. The questionnaire score $V_\beta$ plays the role of $P_\beta$ (alertness), and $V_\alpha$ that of $P_\alpha$ (resting state). Because the questionnaires are administered per session, $E_v$ fixes the overall level of engagement but cannot follow it within a session; that role falls to the EEG decoder.

\paragraph*{Learnable Gramian angular field (LGAF).} To recover the dynamics, short EEG windows are converted into images that a convolutional network can classify. The representation is the learnable Gramian angular field introduced in \cite{HernandezGomez2025GAF} for EEG-based brain--computer interfaces. An EEG window is summarized by the feature matrix
\begin{equation}\label{eq:xeeg}
X_{\mathrm{EEG}}=\begin{pmatrix}V_{\alpha,0}&V_{\alpha,1}&\cdots&V_{\alpha,n}\\ V_{\beta,0}&V_{\beta,1}&\cdots&V_{\beta,n}\end{pmatrix},
\end{equation}
whose two rows are the $\alpha$ and $\beta$ series over the $n+1$ samples of the window, and which the LGAF turns into an image through four successive operations, beginning with a learnable embedding,
\begin{equation}\label{eq:lgaf_embed}
X_t=W_xX_{\mathrm{EEG}}+b_x ,
\end{equation}
whose weights $W_x$ and bias $b_x$ are trained together with the classifier. It then rescales the embedded series to $[-1,1]$,
\begin{equation}\label{eq:lgaf_scale}
\hat X=2\,\frac{x-\min(x)}{\max(x)-\min(x)}-1,
\end{equation}
where, so that gradients can reach $W_x$, the extrema are replaced by their LogSumExp (soft) approximations,
\begin{equation}\label{eq:lse}
\max(x_1,\dots,x_n)\approx\frac1{\alpha_s}\log\sum_{i=1}^{n}e^{\alpha_s x_i},\qquad \min(x_1,\dots,x_n)\approx-\frac1{\alpha_s}\log\sum_{i=1}^{n}e^{-\alpha_s x_i},
\end{equation}
with $\alpha_s>0$ setting the sharpness of the approximation, which becomes exact as $\alpha_s\to\infty$. Each rescaled value is then encoded as an angle,
\begin{equation}\label{eq:arccos}
\phi=\arccos(\hat X)\in[0,\pi],
\end{equation}
and, finally, the angular series of the two rows are compared in a Gramian angular difference field,
\begin{equation}\label{eq:gadf}
\mathrm{GADF}_{ij}=\sin\!\big(W_1\phi^{(1)}_i-W_2\phi^{(2)}_j+b\big),
\end{equation}
in which $\phi^{(1)}$ and $\phi^{(2)}$ are the angular series of the two rows and $W_1$, $W_2$ and $b$ are learnable. Entry $(i,j)$ compares the first series at time $i$ with the second at time $j$, so that the diagonal of the image encodes the instantaneous $\alpha$--$\beta$ relation and its off-diagonal entries encode temporal dependencies between the two bands.

\paragraph*{Decoding cognitive states.} The GADF image is processed by a convolutional network $f_\theta$ with four depthwise-separable convolutional blocks (dilation rates 2 and 4, dropout 0.3 and 0.5), global average pooling and a dense head. The network was trained with AdamW (learning rate $10^{-4}$) on the Nencki--Symfonia dataset \cite{ds004621:1.0.3} to recognise three cognitive states, $c\in\{\text{rest},\text{oddball},\text{MSIT}\}$, corresponding to rest, an oddball task and the Multi-Source Interference Task. On held-out windows it reached an accuracy of 0.71, and the oddball state, which defines the attentive pole of the score, was recognized with the highest F1 (0.74; Table~\ref{tab:lgaf_class}). Applied to a participant's own EEG at time $t$, the network returns logits $z_c(t)$, which the softmax converts into class likelihoods,
\begin{equation}\label{eq:softmax}
L_c(t)=\frac{\exp z_c(t)}{\sum_{c'}\exp z_{c'}(t)} .
\end{equation}

\begin{table}[!htb]
\centering
\caption{\textbf{LGAF-CNN classification on Nencki--Symfonia} (held-out windows): confusion matrix (counts; rows = true class) and per-class metrics. The decoder provides $L_{\mathrm{oddball}}/L_{\mathrm{rest}}$ for CFS.}
\vspace{0.5em}
\begin{tabular}{lccc|cccc}
\toprule
& \multicolumn{3}{c|}{\textbf{Predicted}} & & & & \\
\textbf{True class} & rest & oddball & msit & \textbf{Precision} & \textbf{Recall} & \textbf{F1} & \textbf{Support} \\
\midrule
rest    & \textbf{9,292} & 3,917 & 5,141 & 0.77 & 0.51 & 0.61 & 18,350 \\
oddball & 802 & \textbf{27,153} & 6,079 & 0.69 & 0.80 & 0.74 & 34,034 \\
msit    & 1,916 & 8,009 & \textbf{25,544} & 0.69 & 0.72 & 0.71 & 35,469 \\
\midrule
accuracy & & & & & & 0.71 & 87,853 \\
macro avg & & & & 0.72 & 0.67 & 0.69 & 87,853 \\
weighted avg & & & & 0.71 & 0.71 & 0.70 & 87,853 \\
\bottomrule
\end{tabular}
\label{tab:lgaf_class}
\end{table}

\paragraph*{The Cognitive Flow Score.} The two ingredients are finally combined in a time-resolved score that multiplies the virtualized engagement ratio by the decoded likelihood ratio between an attentive (oddball) state and a resting state,
\begin{equation}\label{eq:cfs}
F(t)=\underbrace{\frac{V_\beta}{V_\alpha}}_{\text{engagement, Eq.~\ref{eq:virtual}}}\cdot\underbrace{\frac{L_{\mathrm{oddball}}(t)}{L_{\mathrm{rest}}(t)}}_{\text{EEG-decoded state}},
\end{equation}
and is sampled at 10\,Hz. The score exceeds $E_v$ whenever the EEG resembles the attentive state more than it resembles rest. Because $E_v$ was nearly constant within a session ($E_v=0.702$ in both conditions), the temporal dynamics of $F(t)$ come from the EEG-decoded ratio, and the comparison between interfaces reported in the Results is, in effect, a comparison of how the decoder reads the two kinds of session.

\paragraph*{LGAF maps of the decision-making network.} Because the decoder operates on whole windows, it is informative to ask what the LGAF representation looks like at the individual nodes of the circuit and whether it changes with engagement. Maps were therefore computed for the nodes of the decision-making circuit (Figure~\ref{fig:circuit}) that the montage covers: core decision and goal maintenance (Fz), motor planning (C3, C4), parietal attention (CP1, CP2, the sites nearest P3/P4) and visual input (O1, Oz, O2); FCz, FC1/2, Fp1/2 and F3/F4 were not recorded. For each node, the raw EEG of each channel was band-passed (fourth-order zero-phase Butterworth) into $\alpha$ (8--12\,Hz) and $\beta$ (13--30\,Hz), and the amplitude envelope was obtained from the analytic signal,
\begin{equation}\label{eq:hilbert}
a_b(t)=\big|x_b(t)+j\,\mathcal H\{x_b\}(t)\big|,\qquad b\in\{\alpha,\beta\},
\end{equation}
where $\mathcal H$ is the Hilbert transform. The envelopes were averaged over the node's channels, interpolated onto the 10\,Hz CFS timeline and cut into windows of $N=64$ steps (6.4\,s), which gave $\approx$\lgafWin{} windows per participant and interface (\lgafWinTotal{} in total). Within each window, each envelope was $z$-scored, rescaled with Eqs.~\eqref{eq:lgaf_scale}--\eqref{eq:lse} ($\alpha_s=4$, clipped to $[-1,1]$) and converted to angles with Eq.~\eqref{eq:arccos}. The map is then the cross-field of Eq.~\eqref{eq:gadf} between the two bands, $\mathrm{GADF}_{ij}=\sin(\phi^{\alpha}_i-\phi^{\beta}_j)$. Because the trained LGAF weights were not archived, the identity initialization ($W_1=W_2=1$, $b=0$) was used. Each map is summarized by its mean diagonal,
\begin{equation}\label{eq:diag}
D=\frac1N\sum_{t=1}^{N}\sin\!\big(\phi^{\alpha}_t-\phi^{\beta}_t\big),
\end{equation}
which is positive when $\beta$ is relatively stronger than $\alpha$. Finally, $D$ was averaged separately over the windows above and below each participant's median CFS, and the two means were compared across participants with a paired $t$-test (Figure~\ref{fig:lgaf}).

\begin{figure}[!htb]
\centering
\figfile{fig_lgaf_maps.pdf}{1.0}
\caption{\textbf{LGAF maps of the decision-making network.} Rows: nodes of the decision-making circuit covered by the montage. Columns: group-averaged maps $\mathrm{GADF}_{ij}=\sin(\phi^{\alpha}_i-\phi^{\beta}_j)$ over 64-step (6.4\,s) windows of the $\alpha$ and $\beta$ envelopes, for windows above (high CFS) and below (low CFS) each participant's median CFS, and their difference, in VR+haptics and 2D. Positive values (red) mean $\beta$ ranks above $\alpha$. Annotations give the high$-$low difference of the diagonal and its paired $p$ value. Maps use the LGAF transform at identity initialization because the trained weights were not archived.}
\label{fig:lgaf}
\end{figure}

\subsection{Behavioral Metrics}
Engagement is only half of the comparison, and the other half, how well the tool was moved, had to be characterized from the shape of the tool trajectory itself, because each block had a fixed duration and completion time could therefore not serve as an outcome. All metrics were computed from the 10\,Hz tool logs of Task~1 and Task~2 and averaged per participant and condition \cite{Varier2022,Chiu2021}, with $p_t\in\mathbb R^3$ denoting the tool position and $\Delta t=0.1$\,s the sampling interval. The four primary metrics capture complementary aspects of control: how abrupt the movement is from one step to the next, how smooth it is over the whole path, how directly each reach heads for its end point, and how much fine deviation rides on top of the intended motion.

\paragraph*{Jerk optimality.} The first metric reuses the optimality term of the original RIDL reward, so that the same quantity describes human performance and later informs the agent. The seven signals $s_{i,t}\in\{y,f_x,f_y,f_z,\mathrm{roll},\mathrm{pitch},\mathrm{yaw}\}$ are differenced and averaged,
\begin{equation}\label{eq:cj}
j_{i,t}=|s_{i,t}-s_{i,t-1}|,\qquad c_t=\frac17\sum_{i=1}^{7}j_{i,t},
\end{equation}
and the result is min--max normalized within the segment and inverted,
\begin{equation}\label{eq:opt}
O_t=1-\frac{c_t-\min_\tau c_\tau}{\max_\tau c_\tau-\min_\tau c_\tau}\in[0,1],
\end{equation}
so that $O_t=1$ at the calmest step of the segment and $O_t=0$ at the most abrupt one.

\paragraph*{Smoothness (LDLJ).} Jerk optimality is relative to each segment, so an absolute measure of smoothness complements it. With velocity $v(t)=\dot p(t)$ and jerk $\dddot p(t)$ estimated by central differences, the log dimensionless jerk of a path of duration $T$ is
\begin{equation}\label{eq:ldlj}
\mathrm{LDLJ}=-\ln\!\left(\frac{T^{3}}{v_{\mathrm{peak}}^{2}}\int_{0}^{T}\big\|\dddot p(t)\big\|^{2}dt\right),\qquad v_{\mathrm{peak}}=\max_t\|v(t)\|.
\end{equation}
The factor $T^3/v_{\mathrm{peak}}^2$ makes the jerk integral dimensionless, and higher (less negative) values indicate smoother movement.

\paragraph*{Path efficiency.} Smoothness says nothing about whether a reach heads straight for its target, and directness is measured separately. The path is smoothed (Savitzky--Golay, 7 samples, order 2) and its speed $u_t=\|\dot p_t\|$ computed. Samples with $u_t>0.2\,\mathrm{median}(u)$ are considered \emph{moving}, and each maximal run of moving samples defines a submovement $k$ with chord-to-arc ratio
\begin{equation}\label{eq:pe}
\eta_k=\frac{\|p_{\mathrm{end},k}-p_{\mathrm{start},k}\|}{\ell_k},\qquad \ell_k=\sum_{t\in k}\|p_{t}-p_{t-1}\|,\qquad \mathrm{PE}=\frac{\sum_k\ell_k\,\eta_k}{\sum_k\ell_k},
\end{equation}
computed over submovements longer than 1\,mm \cite{Jackson2013}, so that perfectly straight reaches give $\mathrm{PE}=1$.

\paragraph*{Micro-deviation (tremor index).} Finally, the fine, tremor-like component of the movement is measured as the RMS distance between the raw path and a $\sim$1\,s Savitzky--Golay smoothing $\tilde p_t$ (11 samples, order 3),
\begin{equation}\label{eq:micro}
\mathrm{MD}=\sqrt{\frac1T\sum_{t=1}^{T}\|p_t-\tilde p_t\|^{2}} .
\end{equation}

Path length $\sum_t\|p_t-p_{t-1}\|$, mean speed, active-time fraction, submovement rate and RMS interaction force were also computed as descriptive measures, whereas grasp success was not logged and is therefore not reported.

\subsection{Statistical Analysis}
Showing that an interface imposes no cost is statistically harder than showing that it helps, because failing to find a difference is not the same as demonstrating its absence. Participant-level comparisons were therefore complemented with equivalence tests and Bayes factors, which can positively support the absence of a meaningful difference. The coupling between engagement and movement over time was analyzed separately on five-second windows, with a mixed-effects model that respects the repeated measurements within each participant.

\paragraph*{Paired tests.} For participant $i$ ($n=10$) and a metric $x$, let $d_i=x_i^{\mathrm{VR}}-x_i^{\mathrm{2D}}$, with mean $\bar d$ and standard deviation $s_d$. The paired $t$ statistic, the effect size and the 90\% confidence interval are
\begin{equation}\label{eq:ttest}
t=\frac{\bar d}{s_d/\sqrt n},\qquad d_z=\frac{\bar d}{s_d},\qquad \mathrm{CI}_{90}=\bar d\pm t_{0.95,\,n-1}\frac{s_d}{\sqrt n},
\end{equation}
with $n-1$ degrees of freedom. Wilcoxon signed-rank tests were run alongside as a check against non-normality, and $p$-values were Holm-corrected within two families, the engagement metrics and the kinematic metrics. With the $m$ ordered $p$-values $p_{(1)}\le\dots\le p_{(m)}$ of a family, the corrected values are
\begin{equation}\label{eq:holm}
\tilde p_{(k)}=\max_{j\le k}\ \min\!\big(1,\,(m-j+1)\,p_{(j)}\big).
\end{equation}

\paragraph*{Equivalence (TOST).} A paired test can only fail to find a difference, so equivalence was tested directly. For an equivalence margin $\pm\Delta$, two one-sided tests are run,
\begin{equation}\label{eq:tost}
t_L=\frac{\bar d+\Delta}{s_d/\sqrt n},\qquad t_U=\frac{\bar d-\Delta}{s_d/\sqrt n},\qquad p_{\mathrm{TOST}}=\max\big\{P(T_{n-1}\ge t_L),\,P(T_{n-1}\le t_U)\big\},
\end{equation}
and equivalence is declared when $p_{\mathrm{TOST}}<0.05$, which is the same as requiring $\mathrm{CI}_{90}\subset(-\Delta,\Delta)$. The primary margin, $\Delta=0.10\,|\bar x^{\mathrm{2D}}|$ ($\pm$10\% of the 2D mean), expresses equivalence in the natural units of each metric; the secondary margin, $\Delta=0.5\,s_d$, corresponds to a medium standardized effect.

\paragraph*{Bayes factors.} Whereas TOST asks whether a difference is confidently small, a Bayes factor weighs the evidence for no difference against the evidence for some difference. Evidence for the null was quantified with the JZS Bayes factor, which places a Cauchy prior with scale $r=\sqrt2/2$ on the standardized effect $\delta$,
\begin{equation}\label{eq:bf}
\mathrm{BF}_{10}=\frac{\int p(t\mid\delta)\,\mathrm{Cauchy}(\delta;0,r)\,d\delta}{p(t\mid\delta=0)},\qquad \mathrm{BF}_{01}=\frac1{\mathrm{BF}_{10}} ,
\end{equation}
and $\mathrm{BF}_{01}>3$ is read as moderate evidence for no difference.

\paragraph*{Mixed-effects model.} Because session averages cannot reveal whether engagement and movement quality rise and fall together within a session, this coupling was modeled on non-overlapping 5\,s windows (50 samples), where for window $j$ of participant $i$
\begin{equation}\label{eq:lmm}
y_{ij}=\beta_0+\beta_1\,c_{ij}+\beta_2\,\mathrm{VR}_{ij}+\beta_3\,\mathrm{Task2}_{ij}+u_{0i}+\varepsilon_{ij},\qquad u_{0i}\sim\mathcal N(0,\tau_0^2),\ \ \varepsilon_{ij}\sim\mathcal N(0,\sigma^2),
\end{equation}
where $y_{ij}$ is the $z$-scored quality metric (optimality, LDLJ or micro-deviation), $c_{ij}$ is the $z$-scored mean CFS of the window, $\mathrm{VR}_{ij}$ and $\mathrm{Task2}_{ij}$ are indicator variables, and $u_{0i}$ is the participant's random intercept. The coefficient $\beta_1$ is the change in quality, in standard deviations, per standard deviation of CFS. Because the strength of such a coupling may differ from person to person, a robustness model adds a random slope, $(\beta_1+u_{1i})\,c_{ij}$ with $(u_{0i},u_{1i})\sim\mathcal N(\mathbf 0,\Sigma)$, so that each participant may have an individual CFS--quality slope. The model was also refitted after 1--99\% trimming and after replacing CFS with its within-participant rank, and all models were fitted by REML. To display the heterogeneity directly, an ordinary least-squares slope was additionally estimated for each participant, with CFS and optimality centred on that participant's block means and scaled by their overall standard deviations.

% -----------------------------------------------------------------
\subsection{RIDL: Hierarchical Reinforcement-and-Imitation Diffusion Learning}
The sessions that measured the learners' engagement also produced demonstrations of how a human moves the surgical tool while engaged to a known degree, and RIDL turns these demonstrations into an agent that predicts the next tool motion from the endoscope image and from the recorded kinematic, force and engagement signals through a two-level hierarchy, in which the agent first imitates the demonstrations and then refines the imitation by reinforcement learning on a reward that weights each moment by the demonstrator's CFS, so that highly engaged behavior counts more (Figure~\ref{fig:ridl}). The description below proceeds from the formal problem to the state and the network that processes it, then to the two levels of learning, and finally to the protocol by which the agent was evaluated.

\begin{figure}[!htb]
\centering
\figfile{fig3_architecture.pdf}{1.0}
\caption{\textbf{Multimodal integration and RIDL architecture.} (1) Three sensing streams: the trainee views the endoscope scene (VR headset or monitor), teleoperates through the Novint Falcon (pose out, force back) and wears 8-channel dry EEG; UWIST/ARSQ questionnaires are collected per session. (2) Per-stream processing: frames are timestamped and resized, kinematics are merged on 0.1\,s timestamps and turned into jerk optimality $O_t$ and displacement history $\Delta a_t$, and EEG passes through the LGAF decoder, whose $L_{\mathrm{oddball}}/L_{\mathrm{rest}}$ ratio is scaled by the virtualized $V_\beta/V_\alpha$ to give CFS$(t)$; all streams are aligned into one 10\,Hz demonstration stream. (3) RIDLPolicy: a frozen Stable-Diffusion UNet encodes the frame, a Conv1D--Transformer encodes the 11-channel history of CFS, $O_t$ and $\Delta a_t$, and the summed 512-d embeddings feed Gaussian policy and value heads. (4) Hierarchical learning: behavior cloning on the operator's 1\,s displacements (level~1) initializes PPO on a reward weighted by CFS (level~2), repeated over participants 1--8; agents are evaluated on participants 9 and 10 with Kalman-filtered residuals. CFS therefore enters twice: as a state input and as the reward weight. The dashed path marks the future closed neuroadaptive loop. Thumbnails show real data from participant~1 (VR, Task~1): endoscope frame and its 64$\times$64 policy input, tool position and force, 8-channel EEG, jerk optimality, an LGAF map, CFS$(t)$ and the four frozen-UNet feature maps of the same frame.}
\label{fig:ridl}
\end{figure}

\paragraph*{Problem formulation.} Robotic teleoperation is modeled as a Markov decision process $(\mathcal S,\mathcal A,P,R,\gamma)$, in which the agent seeks a policy $\pi_\theta(a\mid s)$ that maximizes the expected discounted return,
\begin{equation}\label{eq:mdp}
J(\theta)=\mathbb E_{\pi_\theta}\Big[\sum_{t=0}^{T}\gamma^{t}R_t\Big].
\end{equation}
RIDL obtains $\pi_\theta$ in two levels, imitation of the demonstrations (level~1) followed by reinforcement refinement on a CFS-weighted reward (level~2). Both levels share the same state, network and action, which are defined first.

\paragraph*{State.} At time $t$ the agent observes $s_t=(I_t,N_t)$, where $I_t$ is the current endoscope frame and $N_t\in\mathbb R^{11\times8}$ is the kinematic history over the last eight steps (0.8\,s). The 11 rows of this history are the CFS $F_t$, the jerk optimality $O_t$ (Eq.~\ref{eq:opt}) and the last-step displacements $a_t-a_{t-1}$ of the nine logged quantities $a=(x,y,z,\mathrm{roll},\mathrm{pitch},\mathrm{yaw},f_x,f_y,f_z)$. Every row is standardized, $\tilde n=(n-m)/s$, with means and standard deviations computed on the training participants only, and the first step of a segment is padded with its own value. Engagement thus enters the state directly, and in the ``w/o CFS input'' ablation the $F_t$ row is set to zero.

\paragraph*{Visual encoder.} Because the endoscope frame is the richest of the inputs and the hardest to learn from when demonstrations are few, RIDL does not train a visual encoder but borrows the representation of a large pretrained generative model, Stable Diffusion~v1.5 \cite{rombach2022highresolution}. The frame is resized to $64\times64$ and mapped to $[-1,1]$, and a zero channel is appended to match the four latent channels of the model:
\begin{equation}\label{eq:unet}
x_t=\big[\,2\,\mathrm{Resize}(I_t)-1\,;\ \mathbf 0\,\big]\in\mathbb R^{4\times64\times64},\qquad u_t=\mathrm{UNet}(x_t,\ \tau=0,\ c=\varnothing)\in\mathbb R^{4\times64\times64}.
\end{equation}
With diffusion timestep $\tau=0$ and null text conditioning $c=\varnothing$, the frozen UNet acts as a deterministic feature extractor. Its flattened output (16\,384 dimensions) is standardized and projected to $d=512$,
\begin{equation}\label{eq:fi}
f_I=W_u\,\mathrm{Dropout}_{0.3}\big(\mathrm{vec}(u_t)\big)+b_u\in\mathbb R^{d}.
\end{equation}
Because the UNet is frozen and deterministic, $u_t$ was computed once per frame and cached, which is numerically identical to recomputing it in every forward pass. Two ablations probe this prior: in ``raw pixels'' $\mathrm{vec}(u_t)$ is replaced by the $64\times64\times3$ pixels, and in ``no vision'' $f_I=0$.

\paragraph*{Kinematic encoder.} The kinematic history is a short multivariate sequence, naturally encoded by a Transformer in which each of the eight columns $\tilde n_k$ of $N_t$ is embedded by a $1\times1$ convolution and summed with a sinusoidal positional encoding,
\begin{equation}\label{eq:pe_sin}
h^{(0)}_k=W_e\tilde n_k+b_e+\mathrm{PE}_k,\qquad \mathrm{PE}_{k,2i}=\sin\!\Big(\frac{k}{10000^{2i/d}}\Big),\quad \mathrm{PE}_{k,2i+1}=\cos\!\Big(\frac{k}{10000^{2i/d}}\Big).
\end{equation}
Four Transformer layers follow, each of which applies multi-head self-attention with 8 heads of size $d_h=64$,
\begin{equation}\label{eq:attn}
\mathrm{Attn}(Q,K,V)=\mathrm{softmax}\!\Big(\frac{QK^{\top}}{\sqrt{d_h}}\Big)V,
\end{equation}
together with residual connections, layer normalization and a GELU feed-forward block of width 256 (dropout 0.1),
\begin{equation}\label{eq:tlayer}
h'=\mathrm{LN}\big(h+\mathrm{MHA}(h)\big),\qquad h^{(\ell+1)}=\mathrm{LN}\big(h'+W_2\,\mathrm{GELU}(W_1h'+b_1)+b_2\big).
\end{equation}
The kinematic feature is the mean over the eight output tokens, $f_N=\tfrac18\sum_{k}h^{(4)}_k\in\mathbb R^{d}$.

\paragraph*{Fusion and actor--critic heads.} The visual and kinematic features are fused additively and passed through a shared layer,
\begin{equation}\label{eq:fusion}
f=f_I+f_N,\qquad z=\mathrm{Dropout}_{0.1}\big(\mathrm{GELU}(W_pf+b_p)\big)\in\mathbb R^{256},
\end{equation}
whose output feeds the two heads of an actor--critic, mirroring the policy and value roles assigned to the decision and reward circuits in Figure~\ref{fig:circuit}C: a Gaussian policy (actor) and a value function (critic),
\begin{equation}\label{eq:heads}
\pi_\theta(a\mid s_t)=\mathcal N\!\big(\mu_\theta(s_t),\,\mathrm{diag}\,\sigma^{2}\big),\quad \mu_\theta(s_t)=W_\mu z+b_\mu\in\mathbb R^{9},\quad \sigma=e^{\boldsymbol\ell},\qquad V_\psi(s_t)=w_v^{\top}z+b_v ,
\end{equation}
where $\boldsymbol\ell\in\mathbb R^9$ is a learned, state-independent log standard deviation initialized at $-0.5$.

\paragraph*{Action.} What the policy predicts matters as much as what it observes, and, following action-chunking diffusion policies \cite{chi2023diffusion,kim2024srt}, the action is the displacement of the tool over the next second ($H=10$ steps), standardized with training statistics, and the predicted position follows from it:
\begin{equation}\label{eq:target}
a^*_t=\frac{a_{t+H}-a_t-m_a}{s_a},\qquad \hat p_{t+H}=p_t+\big(s_a\odot\mu_\theta(s_t)+m_a\big)_{x,y,z}.
\end{equation}

\paragraph*{Level 1: behavior cloning.} The first level gives the agent a skill prior by regressing the policy mean onto the demonstrated actions,
\begin{equation}\label{eq:bc}
\mathcal L_{\mathrm{BC}}(\theta)=\frac1{9|\mathcal D|}\sum_{t\in\mathcal D}\big\|\mu_\theta(s_t)-a^*_t\big\|^{2},
\end{equation}
which is minimized with AdamW (learning rate $3\times10^{-4}$, weight decay $10^{-2}$, batch 64, 10 epochs per participant).

\paragraph*{Engagement-weighted reward.} The second level requires a reward, and it is here that engagement enters the learning objective: at each step the agent samples $a_t\sim\pi_\theta(\cdot\mid s_t)$, and its reward is assembled from three terms, the first of which is a tracking kernel, with $d=9$,
\begin{equation}\label{eq:track}
k_t=\exp\!\Big(-\frac{\|a_t-a^*_t\|^{2}}{2\sigma_r^{2}d}\Big)\in(0,1],
\end{equation}
which equals 1 when the agent reproduces the demonstrated chunk and decays with the per-dimension squared error ($\sigma_r=0.5$ in standardized units). The second term, an engagement weight,
\begin{equation}\label{eq:weight}
w_t=\tfrac12\,\mathrm{clip}\big(F_t/\bar F,\,0,\,2\big)\in[0,1],
\end{equation}
in which $\bar F$ is the segment mean, scales the tracking term by the demonstrator's engagement, so that a step taken at average engagement receives $w_t=\tfrac12$, a step at twice the average receives $w_t=1$, and a step at zero engagement receives nothing. The third, a smoothness term,
\begin{equation}\label{eq:smooth}
\varsigma_t=\exp\!\Big(-\frac{\|a_t-a_{t-1}\|^{2}}{2d}\Big),
\end{equation}
penalizes abrupt changes between consecutive actions and is the agent-side analogue of jerk optimality, and together the three terms form the total reward,
\begin{equation}\label{eq:reward}
R_t=A\,w_t\,k_t+B\,\varsigma_t,\qquad A=1,\ B=0.1 ,
\end{equation}
in which the high-engagement moments of a demonstration carry more weight; in the ``w/o CFS-weighted reward'' ablation, $w_t\equiv1$. This reward differs in one important respect from the original RIDL formulation, $R_t=A(1-|F_t-1|)+B\,O_t$, which depended only on logged quantities and was therefore returned whatever the agent did, so that it carried no learning signal, whereas Equation~\eqref{eq:reward} makes the reward depend on the agent's action $a_t$.

\paragraph*{Level 2: PPO refinement.} With the reward defined, the imitation prior is refined by Proximal Policy Optimization. An episode is one demonstration segment, and because the next state is read from the demonstration, actions change the reward but not the state transition, a demonstration-replay setting whose implications are taken up in the Discussion. After a rollout of $T_s$ steps, the temporal-difference residuals are computed with $V_\psi(s_{T_s})=0$,
\begin{equation}\label{eq:td}
\delta_t=R_t+\gamma V_\psi(s_{t+1})-V_\psi(s_t),
\end{equation}
and generalized advantage estimation (GAE) accumulates them backwards in time into advantages and returns,
\begin{equation}\label{eq:gae}
\hat A_t=\sum_{l\ge0}(\gamma\lambda)^{l}\delta_{t+l}=\delta_t+\gamma\lambda\hat A_{t+1},\qquad \hat G_t=\hat A_t+V_\psi(s_t),
\end{equation}
with $\gamma=0.99$ and $\lambda=0.95$. After the advantages have been standardized within the rollout, the policy is updated through the probability ratio between the current parameters and the parameters $\theta_{\mathrm{old}}$ that generated the rollout,
\begin{equation}\label{eq:ratio}
\rho_t(\theta)=\frac{\pi_\theta(a_t\mid s_t)}{\pi_{\theta_{\mathrm{old}}}(a_t\mid s_t)} .
\end{equation}
The clipped surrogate objective built on this ratio keeps the updated policy close to the one that collected the data, and it is accompanied by a value loss and by the entropy of the Gaussian policy,
\begin{equation}\label{eq:clip}
\mathcal L^{\mathrm{CLIP}}=\mathbb E_t\Big[\min\big(\rho_t\hat A_t,\ \mathrm{clip}(\rho_t,1-\epsilon,1+\epsilon)\,\hat A_t\big)\Big],\qquad \mathcal L^{V}=\mathbb E_t\big[(V_\psi(s_t)-\hat G_t)^{2}\big],
\end{equation}
\begin{equation}\label{eq:entropy}
\mathcal H=\sum_{j=1}^{9}\tfrac12\log\!\big(2\pi e\,\sigma_j^{2}\big),
\end{equation}
with $\epsilon=0.2$. The three terms are combined into the loss
\begin{equation}\label{eq:ppo}
\mathcal L_{\mathrm{PPO}}=-\mathcal L^{\mathrm{CLIP}}+c_v\,\mathcal L^{V}-c_e\,\mathcal H,\qquad c_v=0.5,\ c_e=10^{-3},
\end{equation}
which is minimized for 4 epochs over minibatches of 64 with AdamW (learning rate $10^{-4}$) and gradient-norm clipping at 0.5, and six such updates are run per participant \cite{Schulman2017,Larsen2021-gg}.

\paragraph*{Subject curriculum.} The two levels are applied participant by participant, so that the agent accumulates experience from one demonstrator after another: the parameters reached after participant $i$ initialize training on participant $i+1$, $\theta^{(i+1)}_0=\theta^{(i)}$. The ``pooled'' ablation instead trains once on all training participants together, with the same total number of PPO updates.

\paragraph*{Evaluation design.} The agent was evaluated so that no participant ever contributed to both training and testing, by dividing the participants into five subject-wise folds of two held-out participants each and repeating training for every fold with two seeds and with demonstrations from each interface (VR or 2D), which gave 20 runs per variant, each tested on both interfaces of its held-out participants.

\paragraph*{Error metrics.} The basic quantity is the position error of the predicted 1\,s displacement, in millimetres,
\begin{equation}\label{eq:err}
r_t=1000\,\big\|\hat p_{t+H}-p_{t+H}\big\|_2 .
\end{equation}
As in the original RIDL evaluation, $r_t$ is smoothed with a one-dimensional Kalman filter under a random-walk model, $x_t=x_{t-1}+w_t$ and $r_t=x_t+v_t$, with $w_t\sim\mathcal N(0,Q)$ and $v_t\sim\mathcal N(0,R_v)$. With $Q=10^{-5}$, $R_v=0.1$, $\hat x_0=r_1$ and $P_0=1$, each step predicts and then updates,
\begin{equation}\label{eq:kalman}
P^-_t=P_{t-1}+Q,\qquad K_t=\frac{P^-_t}{P^-_t+R_v},\qquad \hat x_t=\hat x_{t-1}+K_t\,(r_t-\hat x_{t-1}),\qquad P_t=(1-K_t)P^-_t ,
\end{equation}
which gives the smoothed residual $\tilde r_t=\hat x_t$, the gain $K_t$ setting how much each new residual moves the estimate. Two summary measures follow from the smoothed residual, the trajectory error and the volatility,
\begin{equation}\label{eq:etraj}
E_{\mathrm{traj}}=\frac1T\sum_{t=1}^{T}\tilde r_t,\qquad \sigma_r=\sqrt{\frac1T\sum_{t=1}^{T}\big(\tilde r_t-E_{\mathrm{traj}}\big)^{2}},
\end{equation}
which describe, respectively, how far and how erratically the prediction deviates from the demonstration (Figure~\ref{fig:kalman}). Because position accounts for only three of the nine predicted quantities, accuracy over all action dimensions is measured separately by the normalized mean squared error,
\begin{equation}\label{eq:nmse}
\mathrm{NMSE}=\frac1{9T}\sum_{t=1}^{T}\|\mu_\theta(s_t)-a^*_t\|^{2},
\end{equation}
computed in standardized units so that position, orientation and force contribute equally.

\paragraph*{Baselines and ablations.} Because a learned policy is only informative relative to what can be achieved without learning, two reference predictors were included: hold position, $\hat p_{t+H}=p_t$, which assumes the tool will not move, and constant velocity, $\hat p_{t+H}=p_t+H\,(p_t-p_{t-1})$, which extrapolates its current motion. Against these references, the ablation removes one component of RIDL at a time: the imitation prior (PPO from scratch), the PPO refinement (BC only), the visual prior (no vision, or raw pixels instead of SD-UNet features), the CFS weighting of the reward ($w_t\equiv1$), CFS as a state input, and the subject curriculum (pooled training). The checkpoints of the original RIDL implementation, one per interface and task, trained sequentially on participants 1--8 with the action-independent reward, were also re-assessed on participants 9 and 10, with targets normalized by training-set statistics so that their errors are comparable with those of the ablation. Finally, sample efficiency was measured by evaluating held-out NMSE after every PPO update, for RIDL and for PPO from scratch, with an extended budget of 15 rollouts per participant.

\section{Results}
Following the order in which the questions were posed, the results show first that immersion left the learners' engagement intact and their movement quality equivalent to that achieved on a conventional display, and that within each session moments of higher engagement coincided with more optimal movement. A closer look at the EEG maps of the decision-making network then shows how this engagement signal is distributed across the circuit. The recordings finally served as demonstrations for the learning agent, whose ablation, sensitivity to the source of demonstrations and sample efficiency close the section.

\subsection{Immersion Leaves Cognitive Engagement Intact}
The immersive interface imposed no detectable cognitive cost on the learner, as participant-level CFS was comparable under VR+haptics ($M=\hCfsVR\pm\hCfsVRsd$) and under the 2D display ($M=\hCfsNVR\pm\hCfsNVRsd$), $t(9)=\hCfsT$, $p\hCfsP$, $d_z=\hCfsDz$ (Figure~\ref{fig:paired}; Table~\ref{tab:subjects}). The 90\% confidence interval of the difference, $[\hCfsCIlo,\ \hCfsCIhi]$, spans $\hCfsCIloPct$\% to $+\hCfsCIhiPct$\% of the 2D mean and therefore excludes any large loss of engagement under immersion, while the Bayes factor mildly favored the absence of a difference (BF$_{01}=\hCfsBF$). Formal equivalence within $\pm$10\% could not be established at this sample size (TOST $p\hCfsTost$), so the result is best read as the absence of a detectable cost rather than as proven equivalence.

If anything, engagement was steadier in the headset, where the within-session coefficient of variation of CFS was \hCfsCvVR{}, compared with \hCfsCvNVR{} under 2D ($p\hCfsCvP$; Holm-corrected $p\hCfsCvPholm$). The 90\% interval of this difference lies below zero, but its 95\% interval does not, and the effect was driven by three participants whose 2D sessions were unusually unstable; it is therefore a trend that a larger sample would need to confirm.

\begin{table}[!htb]
\centering
\caption{\textbf{Participant-level engagement and movement quality.} Means over Task~1 and Task~2. CFS $=(V_\beta/V_\alpha)(L_{\mathrm{oddball}}/L_{\mathrm{rest}})$; path efficiency = submovement chord/arc; $O$ = jerk optimality; LDLJ = log dimensionless jerk (higher = smoother).}
\vspace{0.5em}
\small
\resizebox{\textwidth}{!}{\begin{tabular}{lcccccccc}
\toprule
& \multicolumn{2}{c}{\textbf{CFS}} & \multicolumn{2}{c}{\textbf{Path efficiency}} & \multicolumn{2}{c}{\textbf{Optimality $O$}} & \multicolumn{2}{c}{\textbf{LDLJ}} \\
\cmidrule(lr){2-3}\cmidrule(lr){4-5}\cmidrule(lr){6-7}\cmidrule(lr){8-9}
\textbf{Participant} & VR & 2D & VR & 2D & VR & 2D & VR & 2D \\
\midrule
P1 & 2.35 & 2.09 & 0.49 & 0.27 & 0.917 & 0.927 & -20.24 & -20.32 \\
P2 & 1.63 & 1.70 & 0.35 & 0.28 & 0.951 & 0.957 & -19.86 & -20.28 \\
P3 & 2.38 & 1.61 & 0.29 & 0.11 & 0.933 & 0.941 & -20.20 & -20.83 \\
P4 & 2.17 & 2.52 & 0.19 & 0.12 & 0.928 & 0.902 & -20.52 & -20.50 \\
P5 & 2.31 & 2.34 & 0.41 & 0.54 & 0.916 & 0.942 & -20.88 & -19.15 \\
P6 & 2.58 & 2.03 & 0.69 & 0.37 & 0.986 & 0.896 & -19.49 & -20.49 \\
P7 & 2.09 & 2.09 & 0.29 & 0.14 & 0.924 & 0.938 & -20.43 & -20.65 \\
P8 & 2.06 & 2.14 & 0.45 & 0.36 & 0.972 & 0.976 & -19.99 & -19.45 \\
P9 & 2.28 & 2.31 & 0.13 & 0.18 & 0.923 & 0.892 & -20.49 & -20.95 \\
P10 & 2.23 & 2.20 & 0.32 & 0.25 & 0.936 & 0.950 & -20.37 & -20.40 \\
\midrule
Mean (SD) & 2.21 (0.25) & 2.10 (0.28) & 0.36 (0.16) & 0.26 (0.13) & 0.939 (0.024) & 0.932 (0.028) & -20.25 (0.39) & -20.30 (0.57) \\
\bottomrule
\end{tabular}
}% complete tabular, generated by make_macros.py
\label{tab:subjects}
\end{table}

\begin{figure}[!htb]
\centering
\figfile{fig_human_paired.pdf}{0.9}
\caption{\textbf{Participant-level differences between VR+haptics and 2D/no-haptics.} For each metric, gray dots are individual participants' differences (VR+haptics $-$ 2D) expressed as a percentage of the 2D group mean; the diamond and thick bar give the mean difference and its 90\% confidence interval, which is the interval evaluated by the equivalence test, and the thin whisker gives the 95\% confidence interval, which excludes zero only when the difference is significant at $p<0.05$. The shaded band marks the $\pm$10\% equivalence bounds: a metric is equivalent when its whole interval lies inside the band. Values beyond the axis are pinned to its edge and labelled. The verdict on the right summarizes the paired $t$-test, the Bayes factor and TOST for that metric; participant-level raw values are given in Table~\ref{tab:subjects}.}
\label{fig:paired}
\end{figure}

\subsection{Motor Quality Is Equivalent Across Interfaces}
A cost of immersion might still have appeared in the hand rather than in the EEG, but movement quality told the same story, and here the evidence was strong enough to claim equivalence: jerk optimality was \hOptVR{}~$\pm$~\hOptVRsd{} with VR+haptics and \hOptNVR{}~$\pm$~\hOptNVRsd{} with the 2D display ($p\hOptP$, BF$_{01}=\hOptBF$), and smoothness, measured as the log dimensionless jerk, was \hLdljVR{}~$\pm$~\hLdljVRsd{} and \hLdljNVR{}~$\pm$~\hLdljNVRsd{}, respectively ($p\hLdljP$, BF$_{01}=\hLdljBF$). Both metrics were statistically equivalent within $\pm$10\% of the 2D mean (TOST $p\hOptTost$ and $p\hLdljTost$). Neither micro-deviation (\hMicroVR{} vs.\ \hMicroNVR{}\,mm, $p\hMicroP$) nor path length ($p\hPathLenP$) differed between interfaces. The wide interval of micro-deviation in Figure~\ref{fig:paired} reflects two participants with much lower values in the headset, whereas six of the ten deviated slightly more in VR.

The only kinematic measure that favored immersion was submovement path efficiency, which was higher in the headset (\hPathEffVR{} vs.\ \hPathEffNVR{}, $d_z=\hPathEffDz$, $p\hPathEffP$ uncorrected). Stereoscopic depth cues plausibly support straighter reaches \cite{Fujie2020,Fiorini2022}, but the effect did not survive Holm correction across the kinematic family ($p\hPathEffPholm$) and is reported as exploratory. Representative tool paths, colored by CFS and optimality, are shown in Figure~\ref{fig:cfs} (top), and they anticipate the next question.

\begin{figure}[!htb]
\centering
\figfile{fig4a_rgb_trajectories.pdf}{0.72}\\[2pt]
\figfile{fig4b_cfs_quality.pdf}{0.9}
\caption{\textbf{Engagement--performance relationship.} (Top) Tool paths encoded as RGB $=(1-\hat F,\hat F,\hat O)$ (contrast-stretched between the 2nd and 98th percentiles within each segment): green = high CFS, red = low CFS, bluer = higher jerk optimality; labels give the CFS-weighted reward the demonstration receives at well-separated points. (Bottom left) Jerk optimality across each participant's own CFS range: windows were binned into CFS quintiles within each participant (CFS and optimality centred on the participant's block means), and the line shows the mean $\pm$ 95\% CI across participants, VR and 2D pooled. (Bottom right) Standardized CFS$\rightarrow$optimality slope for each participant ($\pm$ 95\% CI; green = CI above zero) and the pooled random-slope estimate (diamond).}
\label{fig:cfs}
\end{figure}

\subsection{Moment-to-Moment Engagement Tracks Movement Optimality}
Because equivalent averages leave open whether the engagement signal is behaviorally meaningful at all, the analysis next asked whether engagement and execution were coupled in time within each session, and across \mmN{} five-second windows they were, since moments of higher CFS were moments of more optimal, lower-jerk control, with a one-standard-deviation increase in CFS predicting a $\beta=\mmBeta$\,SD increase in jerk optimality (SE $=\mmSE$, $p\mmP$; Figure~\ref{fig:cfs}, bottom). The CFS$\times$interface interaction was not significant ($p\mmInterP$), indicating that the coupling was the same in VR and in 2D.

The strength of the coupling, however, varied from person to person, so that when each participant was allowed an individual slope the average effect remained similar but was no longer significant ($\beta=\mmRsBeta$, $p\mmRsP$), and the same held after trimming the 1\% tails ($\beta=\mmTrimBeta$, $p\mmTrimP$) and on within-subject ranks ($\beta=\mmRankBeta$, $p\mmRankP$). The per-participant slopes make this heterogeneity visible (Figure~\ref{fig:cfs}, bottom right): the slope was positive in \slPos{} of 10 participants, its 95\% confidence interval lay above zero in \slSigPos{}, and it was reliably negative in \slSigNeg{}. In line with this pattern, windows above each person's median CFS showed higher optimality in \msNup{} of 10 participants (\msHi{} vs.\ \msLo{}, $d_z=\msDz$, $p\msP$). The coupling was also specific to optimality, since CFS did not predict smoothness ($\beta=\mmLdljBeta$, $p\mmLdljP$) or micro-deviation ($\beta=\mmMicroBeta$, $p\mmMicroP$), which suggests that engagement relates to how efficiently movements are planned rather than to fine, tremor-like deviations.

\subsection{LGAF Maps of the Decision-Making Network}
If CFS tracks movement optimality, it is natural to ask where in the decision-making circuit the signal originates, and the LGAF maps of each node were inspected for this purpose (Figure~\ref{fig:lgaf}). The group-averaged maps were dominated by negative values in every node and in both interfaces, which means that within most windows the $\alpha$ envelope ranked higher than the $\beta$ envelope. Splitting each participant's windows at the median CFS did not change the $\alpha$--$\beta$ diagonal in any node or interface (all $|d_z|\le\lgafMaxDz$, all $p\ge\lgafMinP$), and the high-minus-low difference maps showed no structure that was consistent across nodes. The engagement captured by CFS therefore does not reduce to a simple $\alpha$/$\beta$ shift in any single node; it arises from the decision boundary that the LGAF-CNN has learned over the whole window.

\subsection{RIDL Ablation}
Having established that the two interfaces yield equivalent behavior and an engagement signal that tracks it, the analysis turned to whether these recordings can teach an agent, beginning with a re-assessment of the checkpoints of the original RIDL implementation, trained sequentially on participants 1--8 with an action-independent reward, on the held-out participants 9 and 10. These checkpoints predicted an almost constant tool pose, with a position error of \tBcErr{}\,mm after imitation and \tRidlErr{}\,mm after PPO fine-tuning, compared with \tBase{}\,mm for always predicting the mean training pose, an outcome that motivated the action-dependent reward and the action-chunk formulation evaluated in the remainder of this subsection.

Table~\ref{tab:ablation} and Figure~\ref{fig:ablation} summarize the cross-validated ablation of the revised agent, with 20 runs per variant (5 subject folds $\times$ 2 demonstration sources $\times$ 2 seeds). Differences are reported as variant minus full RIDL over matched folds, sources and seeds, so that negative values favor the variant, and they were tested with the Wilcoxon signed-rank test.

\begin{table}[!htb]
\centering
\caption{\textbf{RIDL ablation.} Subject-wise 5-fold cross-validation $\times$ 2 seeds $\times$ 2 demonstration sources; errors on held-out participants in both interfaces. Error = 3D position error of the 1\,s displacement (mean $\pm$ 95\% CI over runs); $E_{\mathrm{traj}}$ and $\sigma_r$ = Kalman-smoothed error and volatility (mm); NMSE over all nine action dimensions; steps = PPO environment interactions.}
\vspace{0.5em}
\small
\begin{tabular}{lccccc}
\toprule
\textbf{Variant} & \textbf{Error (mm)} & $\bm{E_{\mathrm{traj}}}$ & $\bm{\sigma_r}$ & \textbf{NMSE} & \textbf{Steps} \\
\midrule
RIDL (full) & 13.79 $\pm$ 3.18 & 13.57 & 2.13 & 4.359 & 14,826 \\
PPO from scratch & 15.33 $\pm$ 1.93 & 15.17 & 1.96 & 4.050 & 14,826 \\
BC only & 10.68 $\pm$ 3.92 & 10.38 & 2.20 & 3.851 & -- \\
w/o SD-UNet (no vision) & 8.42 $\pm$ 1.14 & 8.20 & 2.03 & 2.457 & 14,826 \\
Raw pixels (no SD prior) & 15.00 $\pm$ 2.00 & 14.83 & 2.50 & 3.758 & 14,826 \\
w/o flow-weighted reward & 13.68 $\pm$ 1.53 & 13.46 & 2.08 & 4.322 & 14,826 \\
w/o CFS state input & 13.63 $\pm$ 3.44 & 13.44 & 2.06 & 4.561 & 14,826 \\
Pooled (no curriculum) & 14.05 $\pm$ 1.97 & 13.76 & 2.15 & 3.678 & 14,865 \\
Constant-velocity baseline & 11.85 $\pm$ 2.57 & 11.15 & 6.73 & 152.087 & -- \\
Hold-position baseline & 4.93 $\pm$ 0.98 & 4.54 & 2.37 & 2.093 & -- \\
\bottomrule
\end{tabular}
% complete tabular, generated by make_macros.py
\label{tab:ablation}
\end{table}

\begin{figure}[!htb]
\centering
\figfile{fig_ablation_ablation.pdf}{0.9}\\[4pt]
\figfile{fig5_vis2_VR_fold4.pdf}{0.85}
\caption{\textbf{RIDL ablation.} (Top) Held-out position error per variant (mean $\pm$ 95\% CI over folds, seeds and demonstration sources) and paired differences relative to RIDL. (Bottom) Two held-out recordings (participant~9, 2D and VR+haptics sessions, Task~1): ground-truth tool path (white), PPO from scratch (red) and RIDL (blue), with the per-step error on the right.}
\label{fig:ablation}
\end{figure}

Across the metrics, full RIDL showed a balanced profile and sat in the middle of the learned variants. Its held-out position error (\aRidlErr{}\,mm; Kalman-smoothed $E_{\mathrm{traj}}=\aRidlEtraj$\,mm) was lower on average than that of PPO trained from scratch (\aPpoErr{}\,mm; \dPpo{}\,mm) and of the raw-pixel variant (\aPixErr{}\,mm; \dPix{}\,mm), although neither difference was significant ($p\dPpoP$ and $p\dPixP$). Its residual volatility ($\sigma_r=\aRidlSig$\,mm) was comparable to that of the other learned variants and lower than that of both reference predictors (hold position \aHoldSig{}\,mm, constant velocity \aCVSig{}\,mm). Over all nine action dimensions its normalized error (NMSE \aRidlNmse{}) was about 35 times lower than that of constant-velocity extrapolation (\aCVNmse{}), whose position error (\aCVErr{}\,mm) it matched ($p\dCVP$). The variants that removed CFS from the reward (\dNoFlowR{}\,mm, $p\dNoFlowRP$) or from the state (\dNoCfsIn{}\,mm, $p\dNoCfsInP$), or that pooled participants instead of following the curriculum (\dPooled{}\,mm, $p\dPooledP$), performed on par with full RIDL.

Two variants were more accurate than full RIDL in position error: behavior cloning alone (\aBcErr{}\,mm; \dBc{}\,mm, $p\dBcP$) and the variant without a visual stream (\aNoVisErr{}\,mm; \dNoVis{}\,mm, $p\dNoVisP$). Both are the lower-capacity members of the family, which is what one would expect when training data are scarce, and each fold was trained on only 16 demonstrations (8 participants $\times$ 2 tasks of about 77\,s each). For the same reason, the hold-position predictor, which simply carries the current pose forward, remained the most accurate reference over the 1\,s horizon (\aHoldErr{}\,mm; \dHold{}\,mm relative to RIDL, $p\dHoldP$). Tool motion in these tasks is slow relative to that horizon, and the learned policies had not yet seen enough demonstrations to predict displacements more accurately than zero. The Kalman-smoothed error $E_{\mathrm{traj}}$ ranked all variants in the same order as the raw error, and the error--volatility plane of Figure~\ref{fig:kalman} shows the balance of the full agent between the two.

\begin{figure}[!htb]
\centering
\figfile{fig_kalman_VR_fold4.pdf}{1.0}
\caption{\textbf{Kalman-filtered residuals.} (Left) For one held-out participant in each interface, the raw residual $r_t=\|\hat x_{t+10}-x_{t+10}\|$ (thin) and its Kalman-smoothed version $\tilde r_t$ (bold, $\pm\sigma_r$ band) for RIDL and PPO from scratch; legend values give $E_{\mathrm{traj}}$ and $\sigma_r$. (Right) Error--volatility plane: $E_{\mathrm{traj}}$ vs.\ $\sigma_r$ for every run of every ablation variant and baseline.}
\label{fig:kalman}
\end{figure}

\subsection{Demonstration Source: VR versus 2D}
Because the two interfaces produced equivalent movement quality, demonstrations recorded in either one were expected to train comparable agents. Full RIDL trained on VR+haptics demonstrations reached \srcVR{}\,mm on held-out participants, and RIDL trained on 2D demonstrations reached \srcNVR{}\,mm (difference $\srcDiff{}$\,mm, 90\% CI [\srcCIlo, \srcCIhi], $p\srcP$). The difference was not significant, but equivalence within $\pm$10\% could not be shown (TOST $p\srcTost$), because run-to-run variability exceeded the margin. Agents trained on VR demonstrations had the lowest error on both test interfaces (Table~\ref{tab:crossdomain}), so 2D demonstrations did not transfer better to 2D test data than VR demonstrations did.

\begin{table}[!htb]
\centering
\caption{\textbf{Demonstration source $\times$ test interface} for full RIDL. Held-out position error (mm, mean $\pm$ SD over folds, seeds and recordings).}
\vspace{0.5em}
\begin{tabular}{lcc}
\toprule
\textbf{Trained on} & \textbf{Tested on VR+haptics} & \textbf{Tested on 2D} \\
\midrule
VR+haptics demonstrations & 11.25 $\pm$ 4.04 & 13.13 $\pm$ 7.22 \\
2D demonstrations & 15.20 $\pm$ 11.23 & 15.58 $\pm$ 11.99 \\
\bottomrule
\end{tabular}
% complete tabular, generated by figures_rl.py
\label{tab:crossdomain}
\end{table}

\subsection{Sample Efficiency}
An imitation prior is expected to make reinforcement learning more sample-efficient, and the last analysis tested this expectation. Over \cSteps{} PPO environment steps (VR demonstrations, 5 folds $\times$ 2 seeds), neither agent improved markedly on its starting point. RIDL started from the imitation prior at a held-out NMSE of \cRidlFirst{} after the first update and ended at \cRidlFinal{}. Its learning curve jumped each time the curriculum moved on to the next participant (about every 4,600 steps), showing that PPO partly overwrote the prior instead of refining it. PPO from scratch started lowest (\cPpoFirst{}), because a freshly initialized Gaussian head predicts a near-zero displacement and thus behaves like the hold predictor, and it then degraded steadily to \cPpoFinal{}. RIDL's curve fluctuated around that of PPO from scratch and ended below it (\cRidlFinal{} vs.\ \cPpoFinal{}), so the imitation prior kept the policy in a better region than random initialization did, although a clear sample-efficiency advantage will only become measurable with more demonstrations per participant.

% =================================================================
\section{Discussion}

This study asked whether an immersive, force-rendering interface for surgical teleoperation changes how learners engage, and whether the recordings such an interface produces can teach a hierarchical learning agent. The answer to the first question is reassuring, the answer to the second is encouraging but incomplete, and the two are linked by the finding that both interfaces produce equivalent behavior.

In this pilot, stereoscopic VR with force feedback did not lower EEG-derived engagement or movement quality relative to a 2D display. The confidence interval on CFS excludes large losses, the Bayes factors mildly favored the null, and movement optimality and smoothness were formally equivalent. These results are interpreted as evidence that the immersive interface is \emph{cognitively transparent}: it adds presence and depth cues without imposing a measurable engagement or motor cost. They do not show that immersion \emph{improves} performance, and no such claim is made here. Cognitive-load theory predicts that richer multisensory interfaces can add extraneous load \cite{jining2021lost,Sokolowska2023-cz}, and the absence of such a cost in this study is encouraging for the use of VR and haptics in surgical curricula \cite{Mao2021,Gawecki2020,9262701}. Transparency matters in a second way for neuroadaptive tutoring \cite{Battaglia2021,Bourdillon2023}, since an engagement signal that is stable across interfaces can be calibrated once and reused whether a trainee works in a headset or at a console.

The engagement signal was also behaviorally meaningful, since windows of higher CFS were windows of lower-jerk, more optimal control, and the coupling was the same in both interfaces. This is consistent with the decision-circuit view adopted in the framework, in which frontal goal maintenance and parietal attention modulate motor planning and output (Figure~\ref{fig:circuit}). The strength of the coupling, however, varied across individuals, in line with trait-level differences in how attention maps onto motor control \cite{Medaglia2020}, and this heterogeneity argues for per-learner calibration of any closed-loop policy that acts on CFS, for example one that slows tool dynamics or adds haptic guidance when engagement drops and raises task difficulty when engagement is high \cite{Attanasio2021,Tagliabue2023,Fiorini2022}. The LGAF maps add a further nuance, because the engagement signal did not appear as a simple $\alpha$/$\beta$ shift in any single node, which suggests that CFS draws on a pattern distributed across the network rather than on a single cortical marker, and that the decoder, not the raw spectral balance, is what carries the moment-to-moment information.

The RIDL ablation gives a mixed but informative picture, in which full RIDL showed a balanced profile across error, volatility and normalized error, outperforming PPO trained from scratch and the raw-pixel variant on average and producing far more stable predictions than constant-velocity extrapolation. The imitation prior therefore did its job of keeping reinforcement learning in a sensible region of policy space \cite{Barnoy2021,Xu2021}. At the same time, the lower-capacity variants, behavior cloning alone and the policy without a visual stream, were more accurate, and none of the learned policies yet improved on carrying the current pose forward.

This pattern is attributed mainly to the amount of data, since each fold was trained on only 16 short demonstrations, a small set for a policy with a 512-dimensional Stable-Diffusion embedding and a four-layer Transformer, whereas diffusion-based visuomotor policies are typically trained on far larger demonstration sets \cite{chi2023diffusion,kim2024srt}. A second factor is the demonstration-replay setting, in which the agent's action changes the reward but not the next state, which limits what PPO can add and also explains why the CFS reward weight and the CFS input made no detectable difference at this stage: an engagement weight can only shape learning once the agent's actions have consequences that differ between engaged and disengaged moments. Neither observation argues against the hierarchical design or against neuroadaptive rewards, but together they indicate that the next steps are more demonstrations per participant, more participants, and a closed-loop simulator in which actions move the instruments, such as the dVRK in AMBF \cite{Munawar2021}. Agents trained on VR and on 2D demonstrations did not differ, consistent with the equivalent motor quality of the two interfaces, which suggests that either interface can supply demonstrations as the data set grows.

These results also sit within a broader trajectory in which brain--computer interfaces and AI decision support point toward higher surgical autonomy under surgeon oversight \cite{Guger2009Brain,Hochberg2006Neuronal,long2022,qian2019}, and a training environment that reads the learner's engagement and learns from the learner's movements is one step along that path. The ethical implications of decision-making autonomy and of biometric monitoring during training must be addressed alongside the technology itself \cite{Medaglia2020,Fosch-Villaronga2023,randell2015impact}.

\subsection{Limitations}
These conclusions are qualified first by the small sample ($N=10$), which is why equivalence for CFS could not be established and why the exploratory path-efficiency effect and the steadier engagement in the headset require replication. They are further qualified by the fixed condition order, because all participants performed VR before 2D, so that practice effects, which would favor 2D, and fatigue, which would favor VR, are confounded with the interface and only a counterbalanced design can separate them. On the measurement side, the virtualized $V_\beta/V_\alpha$ term is defined per session, so the dynamics of CFS rely entirely on the LGAF-CNN decoder, which was trained on a different cohort and task set (Nencki--Symfonia, 71\% accuracy) and transferred to 8-channel dry EEG recorded during movement \cite{Duvinage2013Performance}. Presence and body ownership were not measured with validated questionnaires, so the study cannot relate engagement to the subjective sense of embodiment. On the agent side, RIDL was trained on only 16 demonstrations per fold, which limits what a high-capacity visuomotor policy can learn, and it was trained and tested on logged demonstrations in simulation, so closed-loop deployment and transfer to physical dVRK hardware remain untested \cite{d2021accelerating,Matas2018}. Finally, complementary biosignals such as galvanic skin response, heart-rate variability and fNIRS could strengthen engagement estimation \cite{Ban2022KernelFlow,Li2024Reduced,Melcer2018CTRLLabs}.

These limitations also outline the next study, in which larger, counterbalanced cohorts with validated presence measures would establish equivalence on firmer ground, closed-loop simulation would allow CFS-driven adaptation and engagement-weighted rewards to act on the agent's own consequences, and long-term retention and transfer studies would test whether neuroadaptive training changes what learners carry into the operating room \cite{Gawecki2020}.

\section{Conclusion}
This study shows that a VR--haptic surgical-teleoperation interface preserves EEG-derived engagement and movement quality relative to a conventional display, and that moment-to-moment engagement tracks movement optimality in both. Immersive interfaces can therefore serve as the substrate for neuroadaptive, learner-aware training \cite{9262701,Feussner2017}, with an engagement signal that does not depend on how the scene is displayed. Used as demonstrations, the same recordings trained a RIDL agent with a balanced error and volatility profile that outperformed reinforcement learning from scratch on average. With only 16 demonstrations per fold, simpler policies and a hold-position reference were still more accurate, and larger demonstration sets together with closed-loop simulation are the natural next step toward realizing the benefit of the imitation prior and of engagement-weighted rewards.

% =================================================================
\section*{Ethics Statement}
All participants provided written informed consent before taking part. The study was non-clinical and involved only non-invasive EEG recording and interaction with a haptic VR simulator.

\section*{Author Contributions}
N.N.G.: conceptualization, system design and integration (VR, haptics, EEG, dVRK simulation), methodology, data collection, model development (CFS, RIDL), formal analysis and writing (original draft). A.H.G.: methodology and software for the LGAF pipeline, including its formulation and the training of the LGAF-CNN, building on the homologous work in \cite{HernandezGomez2025GAF}. G.D.L.R.A.: software and resources for the VR integration. F.J.F.A.: methodology, design of the agent that learns from the demonstrations within RIDL. A.J.M.L.: investigation, participant testing and data collection. L.A.M.U.: project administration, supervision and funding acquisition. All authors contributed to writing (review and editing) and approved the submitted version.

\section*{Acknowledgments}
The authors thank the IFE Living Lab for access to immersive simulation facilities and haptic hardware, and the volunteer participants. Andr\'es Ayala at Mostla provided essential devices and technical guidance.

\section*{Data and Code Availability}
Code for synchronized VR/haptics/EEG logging, preprocessing, CFS derivation, the RIDL implementation and all analyses in this article (including the scripts that generate every reported number) will be released upon publication at \url{https://github.com/nezih-niegu/neuroadaptive-ridl}. Anonymized derived features (CFS time series and kinematic metrics) will accompany the release; raw EEG and identifiable data will not be released.

\section*{Conflict of Interest}
The authors declare no commercial or financial relationships that could be construed as a potential conflict of interest.

% =================================================================
\clearpage

\bibliographystyle{plain}
\bibliography{test}
\end{document}
